

 title\+: External Communication permalink\+: doc\+\_\+external\+Comm.\+html keywords\+: design sidebar\+: doc\+\_\+sidebar \subsection*{folder\+: doc }

The main concept of the external\+Comm library is similar to the F\+MI Standard and it can be viewed as blackbox with time-\/dependent input and output values and parameters that are fixed in time.

\section*{Usage}

\subsection*{add to Open\+F\+O\+AM}

add following library to the control\+Dict


\begin{DoxyCode}
libs(externalComm)
\end{DoxyCode}


\subsection*{input}


\begin{DoxyCode}
\textcolor{comment}{// create new entry}
\textcolor{keyword}{const} Time& runTime = mesh\_.time(); \textcolor{comment}{// we need the singleton time}
commDataLayer& data = commDataLayer::New(runTime);

scalar varName = \textcolor{stringliteral}{"varName"};
scalar varValue = \textcolor{stringliteral}{"varValue"};

data.storeObj(varValue,varName,commDataLayer::causality::in);
\end{DoxyCode}



\begin{DoxyCode}
\textcolor{comment}{// work with data}
commDataLayer& data = commDataLayer::New(runTime);
\textcolor{comment}{// get varName from registry}
scalar& val = data.getObj<scalar>(\textcolor{stringliteral}{"varName"},commDataLayer::causality::in);
val = newValue;
\end{DoxyCode}


\subsection*{output}


\begin{DoxyCode}
\textcolor{comment}{// create new entry}
\textcolor{keyword}{const} Time& runTime = mesh\_.time(); \textcolor{comment}{// we need the singleton time}
commDataLayer& data = commDataLayer::New(runTime);

scalar varName2 = \textcolor{stringliteral}{"varName2"}; 

vector varValue2(0,0,0); 
data.storeObj(varValue2,varName2,\hyperlink{createOutput_8H_a0467985694a37c29b5874ecb748bf09e}{commDataLayer::causality::out});
\end{DoxyCode}



\begin{DoxyCode}
\textcolor{comment}{// work with data}
commDataLayer& data = commDataLayer::New(runTime);
\textcolor{comment}{// get varName2 from registry}
vector& val2 = data.getObj<vector>(\textcolor{stringliteral}{"varName2"},\hyperlink{createOutput_8H_a0467985694a37c29b5874ecb748bf09e}{commDataLayer::causality::out}); 
val2 = newValue;
\end{DoxyCode}


\section*{Implementation}

\subsection*{Data}

As hinted above the approach is similar to F\+MI standard with the categorization of variables in 3 causality\+:


\begin{DoxyItemize}
\item input -\/$>$ changes in time (and space)
\item output -\/$>$ changes in time (and space)
\item parameter -\/$>$ constant in time (and space)
\end{DoxyItemize}

The core class {\ttfamily comm\+Data\+Layer} creates 3 {\ttfamily object\+Registries} one for each causality that are accessible with an enum\+:


\begin{DoxyItemize}
\item {\ttfamily comm\+Data\+Layer\+::causality\+::in}
\item {\ttfamily comm\+Data\+Layer\+::causality\+::parameter}
\item {\ttfamily comm\+Data\+Layer\+::causality\+::out}
\end{DoxyItemize}

Each of these {\ttfamily object\+Registries} store and manage the objects passed to them. The elements in the objects\+Registry can be accessed similar to a hashtable. The data is stored in the template class {\ttfamily registered\+Object$<$Type$>$} which currently features following Types\+:


\begin{DoxyItemize}
\item Primitive Types\+:
\begin{DoxyItemize}
\item {\ttfamily bool} -\/$>$ F\+MI 2.\+0
\item {\ttfamily word} -\/$>$ F\+MI 2.\+0
\item {\ttfamily label} -\/$>$ F\+MI 2.\+0
\item {\ttfamily scalar} -\/$>$ F\+MI 2.\+0
\item {\ttfamily vector}
\item {\ttfamily symm\+Tensor}
\item {\ttfamily spherical\+Tensor}
\item {\ttfamily tensor}
\end{DoxyItemize}
\item Fields\+:
\begin{DoxyItemize}
\item {\ttfamily Field$<$scalar$>$}
\item {\ttfamily Field$<$vector$>$}
\item {\ttfamily Field$<$symm\+Tensor$>$}
\item {\ttfamily Field$<$spherical\+Tensor$>$}
\item {\ttfamily Field$<$tensor$>$}
\end{DoxyItemize}
\end{DoxyItemize}

\subsection*{Comm}

\subsubsection*{Websocket}

\subsubsection*{Zmq}

\subsubsection*{grpc}

\subsection*{Json}

\subsection*{example}